\section{转动、振动和量子涨落}\label{sec:3-1}

在第2.4节与第2.6节中，我们建立了非轴对称原子核的玻尔四极形状坐标、主轴参考系与积分测度，剖析了势能曲面刚度、转动惯量与微观配对响应的解耦机制；而在后续第3.2节中，将进一步考察固定惯量下三轴转子绕主惯量轴的小振幅局域微扰展开。然而，面对真实实验所观测到的复杂在束能谱与电磁跃迁网络，初学者极易产生两类认识误区：一是将能谱上的自旋宇称标签误同化为经典刚体粒子在实空间中的确定几何轨道；二是把势能曲面上的静态极小点误当成原子核的实际物理构型，忽略了有限介观量子波包具有不可忽视的零点运动涨落。

作为第3章探讨三轴与多极激发动力学的总览入口，本节旨在系统厘清三个不同维度的物理概念：大范围空间取向转动、局部内禀形状振动与有限波包的量子零点涨落。我们从一维教学振子模型出发，明确恢复力常数与质量参数的物理量纲，解析零点涨落宽度的本质；继而探讨三维空间转动量子化、Casimir 能标与基态电四极矩的谱学不变量；进而阐明内禀形状坐标与欧拉取向角的分离条件及声子数标签的近似边界；深入剖析浅软势能面、多阱共存与量子隧穿在有限体系中的表象；最后平滑引出后续章节的动力学演化脉络。

\subsection{一维教学振子：恢复力、质量参数与零点涨落}\label{sec:3-1-1}

为了给多体坐标下的复杂量子涨落建立坚实的物理参照，首先考察一个沿广义标量坐标 $q$ 运动的一维教学简谐振子。设体系动能为 $T = p^2 / (2B)$，势能在局部极小点 $q = 0$ 附近作二次泰勒微扰展开：
\begin{equation}\label{eq:3-1-1d-potential}
V(q) = V_0 + \frac{1}{2} C q^2 \quad (C > 0),
\end{equation}
其中势能极小处于稳定平衡点，要求刚度系数满足 $C > 0$。

在物理量纲上，需要严格规范广义坐标与动力学参数的单位对应关系：
\begin{enumerate}
\item 广义坐标 $q$ 可以取不同的物理对象：若 $q$ 代表内禀四极形变度 $\alpha$ 或 $\beta$，则 $q$ 为无量纲量；若 $q$ 代表核表面的径向伸缩位移，则其量纲为长度（$[q] = \mathrm{fm}$）；若 $q$ 代表多极跃迁矩，则其量纲为电荷乘面积（$[q] = e\cdot\mathrm{fm}^2$）；
\item 质量参数 $B$ 决定了广义动量 $p = B \dot{q}$ 与动能 $T = p^2 / (2B)$ 的对应关系。为了使动能具备能量量纲 $\mathrm{MeV}$，质量参数 $B$ 的单位严格为 $\hbar^2 / (\mathrm{MeV} \cdot [q]^2)$（若 $q$ 无量纲，则 $B$ 的单位为 $\hbar^2 / \mathrm{MeV}$）；
\item 刚度系数（恢复力常数）$C \equiv \left( \frac{\partial^2 V}{\partial q^2} \right)_{q=0}$ 决定了势阱曲率，其物理单位严格为 $\mathrm{MeV} / [q]^2$。
\end{enumerate}

对该单自由度哈密顿量进行标准正则量子化，基本对易关系为 $[\hat{q}, \hat{p}] = i\hbar$。体系特征谐振角频率（能量标度）由质量与刚度的比值严格给出：
\begin{equation}\label{eq:3-1-oscillator-freq}
\hbar \omega = \hbar \sqrt{\frac{C}{B}} \quad (\mathrm{MeV}).
\end{equation}
基态量子波函数是定域于势阱底部的实高斯波包：
\begin{equation}\label{eq:3-1-ground-state-wf}
\psi_0(q) = \left( \frac{B \omega}{\pi \hbar} \right)^{1/4} \exp\left( -\frac{B \omega}{2\hbar} q^2 \right).
\end{equation}
基态量子期望值直接给出广义坐标的零点涨落均方根宽度：
\begin{equation}\label{eq:3-1-zero-point-fluc}
\langle q^2 \rangle_0 = \int_{-\infty}^{+\infty} |\psi_0(q)|^2 q^2 \, dq = \frac{\hbar}{2 B \omega} = \frac{\hbar}{2 \sqrt{B C}}.
\end{equation}

式~\eqref{eq:3-1-zero-point-fluc} 给出基态波包的坐标方差；其平方根就是均方根零点涨落宽度。它由所选势阱的质量与恢复力决定，是量子态自身的统计矩。不确定性原理在此仅给出乘积下界：
\begin{equation*}
\Delta q \Delta p \ge \frac{\hbar}{2} .
\end{equation*}
在物理理解与实验测量之间，需要系统澄清波包统计矩、跃迁矩阵元与测量误差的辩证关系：
\begin{itemize}
\item \textbf{波包统计矩与非对角跃迁矩阵元的完备性联系}：对于线性坐标算符 $\hat{q}$，当体系具有良好基底完备性且基态平均值 $\langle 0 | \hat{q} | 0 \rangle = 0$ 时，利用完备性求和（closure 关系）可得：
\begin{equation}\label{eq:3-1-closure-q2}
\langle 0 | \hat{q}^2 | 0 \rangle = \sum_n |\langle n | \hat{q} | 0 \rangle|^2.
\end{equation}
这表明基态波包涨落平方并非与可观测量完全割裂，它严格对应于从基态出发的跃迁响应网络强度总和。但这并不意味着可将单个非对角跃迁矩阵元 $\langle f | \hat{\mathcal{M}} | i \rangle$（如低激发态间的 $B(E2)$ 跃迁）直接等同于基态涨落宽度本身；
\item \textbf{量子波包宽度与经典静态极小点的本质区别}：经典力学中体系可静态定域于势能谷底 $q_0 = 0$ 且无动能涨落；在量子力学中，即使势能在原点具有极深势阱，零点振动能与基态波包弥散度 $\sqrt{\langle q^2 \rangle_0} > 0$ 始终存在。波包的统计涨落宽度与实验仪器的多普勒测量误差或谱线展宽亦分属不同范畴。
\end{itemize}

\subsection{转动量子化、Casimir 标度与谱学不变量}\label{sec:3-1-2}

当我们将上述一维几何扩展到三维多体转动时，体系呈现出全新的对称性守恒律与量子化阶梯。孤立原子核在无外场空间中满足空间各向同性，其总角动量守恒且具有固定的 Casimir 模长 $\hat{\mathbf{J}}^2 = \hbar^2 I(I+1)$。

三维刚性转子的哈密顿量写作：
\begin{equation}\label{eq:3-1-rot-H}
\hat{H}_{\text{rot}} = \sum_{k=1}^3 \frac{\hat{J}_k^2}{2\mathcal{J}_k} = \sum_{k=1}^3 A_k \mathcal{I}_k^2,
\end{equation}
其中 $\mathcal{I}_k \equiv \hat{J}_k / \hbar$ 为无量纲体轴角动量算符，$A_k \equiv \frac{\hbar^2}{2\mathcal{J}_k}$ 为能量量纲转动常数（$\mathrm{MeV}$ 或 $\mathrm{keV}$），与第3.2节中的约定严格保持一致。在具有轴对称动力学惯量 $\mathcal{J}_1 = \mathcal{J}_2 \equiv \mathcal{J}_\perp$ 且内禀投影 $K=0$ 的转动带中，转动能谱表述为：
\begin{equation}\label{eq:3-1-rot-spectrum}
E_{\text{rot}}(I) = \frac{\hbar^2}{2\mathcal{J}_\perp} I(I+1) = A_\perp I(I+1).
\end{equation}
需要明确的是，$K=0$ 条件本身仅要求角动量沿对称轴的分量为零，并不迫使自旋 $I$ 只能取偶数（正如第2.5节所严格推导的，$I$ 可取任意非负整数，每个自旋包含 $(2I+1)$ 重磁简并 $M = -I, \dots, +I$）。基带中仅出现偶数自旋序列 $I = 0, 2, 4, \dots$ 的物理前提，是内禀参考态额外具备绕垂直于对称轴的主轴旋转 $\pi$ 的离散二度轴对称性（正旋称 $R_{y, \text{intr}}(\pi) |\Phi\rangle = + |\Phi\rangle$）；若不具备该内禀离散对称性，奇数自旋同样可以存在。

在此必须厘清角动量投影量子数与模型假设的依赖关系：
\begin{enumerate}
\item \textbf{内禀投影 $K$ 的模型属性}：在理想轴对称强变形核中，自旋在内禀对称轴上的投影 $K$ 是一个优良的绝热近似量子数。然而在一般的三轴非轴对称核中（$A_1 \ne A_2 \ne A_3$），哈密顿量打破了绕单一主轴旋转的连续对称性，不同 $K$ 分量发生严重混合，$K$ 不再是守恒标签，物理态是由多个不同 $K$ 态构成的相干叠加（详见第2.5节与第3.2节）；
\item \textbf{基态静态四极矩 $Q_s(0^+) = 0$ 绝不否定核形变}：在实验谱学中，偶偶核基态的自旋宇称恒为 $I^\pi = 0^+$。根据第1.2节阐释的 Wigner-Eckart 定理，任意处于自旋为 $I$ 的态矢对二阶球张量电四极算符 $\hat{Q}_{20}$ 的对角期望值正比于 Clebsch-Gordan 系数 $(I \, 2 \, I; M \, 0 \, M)$。当自旋 $I = 0$ 时：
\begin{equation}\label{eq:3-1-WE-zero}
(0 \, 2 \, 0; 0 \, 0 \, 0) \equiv 0 \implies Q_s(0^+) = \sqrt{\frac{16\pi}{5}} \langle 0, 0 | \hat{Q}_{20} | 0, 0 \rangle = 0.
\end{equation}
式~\eqref{eq:3-1-WE-zero} 是全空间三维角动量几何耦合的必然数学结论，与原子核内禀是否发生四极形变毫无关联。即使核在内禀系具备极大的四极变形，其实验室基态各向同性取向平均后测得的静态四极矩依然严格为零。真实的内禀电四极度由实验室系转动标量不变量所确定（Kumar-Cline 形状不变量，Cline 1986，定义与第2.1节严格统一）~\cite{cline-1986-nuclear-shapes-coulomb-excitation}：
\begin{equation}\label{eq:3-1-cline-invariant}
q_2 \equiv \sqrt{5} \langle 0^+ | [\hat{Q} \times \hat{Q}]_0 | 0^+ \rangle = \sum_i |\langle i || \hat{\mathcal{M}}(E2) || 0^+ \rangle|^2 \ge 0.
\end{equation}
这一表达式使用完整中间态求和。$q_2$ 刻画所用电四极算符的均方强度；将它进一步解释为内禀振幅或形状分布，要保留第~2.1、2.4~节的条件。单个矩或少数跃迁不能代替完整分布。
\end{enumerate}

\subsection{形状振动与取向运动的坐标解耦及声子标签界限}\label{sec:3-1-3}

低能原子核结构理论在描述集体动力学时，将五维空间中的多极自由度 $\alpha_{2\mu}$（$\mu = -2, \dots, +2$）解耦为两套物理本质截然不同的运动坐标（Bohr 1952）~\cite{bohr-1952-surface-particle-coupling}：
\begin{itemize}
\item 2 个内禀几何形状变量 $(\beta, \gamma)$：刻画核椭球的径向四极拉长与三轴非轴对称横向剪切，对应核的\textbf{内禀形状振动}；
\item 3 个欧拉取向角 $\Omega = (\phi, \theta, \psi)$：确定内禀主轴坐标系相对于实验室空间的欧拉转角，对应核的\textbf{大范围整体转动}。
\end{itemize}

对于所选参考态，转动能标小于相关内禀激发能隙，是采用振转绝热近似的常用条件之一\cite{book-bohr-mottelson-nuclear-structure-vol2}：
\begin{equation}\label{eq:3-1-adiabatic-condition}
\hbar \omega_{\text{rot}} \ll \hbar \omega_{\text{vib}}.
\end{equation}
需要澄清的是，绝热分离的本质是转动时间尺度远慢于内禀响应时间尺度，由转动角速度诱导的非绝热导数耦合（如科里奥利微扰 $\langle \Phi_1 | \frac{\partial}{\partial \Omega} | \Phi_2 \rangle$）较小；这与波包几何振幅的大小（$\Delta R / R_0$）属于不同物理层次的评估。大振幅集体运动在能隙充足时依然可以保持绝热演化，而微小振幅在能级微观反交叉（pseudo-crossing）或高能级密度区域亦可诱发剧烈的非绝热混合。

在绝热近似下，转动与内禀振动波函数可表示为对称化的乘积试探形式：
\begin{equation}\label{eq:3-1-factorized-wf}
\Psi_{IM}(\beta,\gamma,\Omega)\approx\sum_K\Phi_K(\beta,\gamma)\,\varphi_{IMK}(\Omega).
\end{equation}
其中 $\varphi_{IMK}$ 是按所用Euler角、离散对称性和测度构造的转动函数；其Wigner $D$ 组合与规范化见第~2.5~节及附录B。本式只示意形状振幅与取向函数的分工，少数 $K$ 项的保留属于具体振转近似，不能视为一般三轴核的精确本征态。

基于上述分离，唯象集体模型常引入简谐声子量子数（如 $\beta$ 振动声子数 $n_\beta$、$\gamma$ 振动声子数 $n_\gamma$ 或摆动声子数 $n_w$）来标记不同能带。然而，必须清醒认识到声子数 $n$ 仅仅是局部简谐展开下的数学标签：
\begin{enumerate}
\item \textbf{势能面的强非简谐性}：真实核势能面包含高阶项（如 $\beta^4$、$\beta^3 \cos(3\gamma)$），微观质量参数与转动惯量亦强烈依赖于动态坐标 $(\beta, \gamma)$，破坏了严格简谐二次型；
\item \textbf{转动—振动耦合效应}：当转动频率增加时，离心力拉伸形变（离心拉伸效应），科里奥利力引发内部准粒子顺排，诱导转动态与振动态发生剧烈动力学杂化；
\item \textbf{整数能量比绝非模式判据}：在束谱学中，实验测得的激发能级比值（如 $E(4_1^+)/E(2_1^+) \approx 2.0$ 对应简谐振动，$E(4_1^+)/E(2_1^+) \approx 3.33$ 对应刚性转子，比值接近 $2.5$ 对应软核或 $\gamma$ 不稳定核）仅是粗糙的唯象线索，绝不能充当认定激发模式的充要判据。判定应结合具体模型框架，考察带间绝对跃迁矩阵元 $B(E2)$、跃迁选择定则以及分支比等综合指标，而不强求每一个具体核都必须测齐全部绝对跃迁强度。
\end{enumerate}

\subsection{浅软势能面、多阱隧穿与量子波包分布}\label{sec:3-1-4}

上述关于简谐振动与定常转动的图景，建立在势能曲面具有显著深势阱极小的假说之上。然而在过渡核区与软核区（如 $A \approx 130$ 核区的部分核素），这一前提往往发生动摇。

\textbf{浅软势能面与大振幅零点涨落}：在许多低能原子核中，微观约束密度泛函（CDFT）或推转微观模型计算给出的势能面沿 $\gamma$ 方向非常平缓（如 Wilets-Jean $\gamma$ 不稳定软势），或沿 $\beta$ 方向势垒极低。此时零点涨落宽度（如 $\sqrt{\langle \Delta \gamma^2 \rangle_0}$）与局部二次谐振展开的有效适用区间相当甚至超界（特别是在长轴对称平衡点 $\gamma_0 = 0^\circ$ 附近，单纯将涨落与零平衡值比较失去意义，必须与势能曲面曲率主导区开展比较）。当零点涨落横跨大范围变形区域时，势能面上的静态极小点不能代表体系的实际有效几何构型；真实物理特性完全由量子波包概率密度分布 $|\Psi(\beta, \gamma)|^2$ 在全流形上的加权分布所主导。

\textbf{多阱共存与量子隧穿}：当微观能量曲面存在多个深度相近的分立极小（例如长椭球极小与扁椭球极小共存，或球形基态与变形异能态共存）时，体系在有限势垒之间发生量子隧穿。在此情形下，体系的能量本征态不再局域于单一势阱，而是分布在两阱中的波包相干叠加形成的宏观量子干涉态（形状共存态）。

\textbf{有限体系定态与内禀对称性破缺的范畴界限}：在宏观固体或液晶中，热力学极限下无穷大自由度的集体冻结会导致实验室系出现真正的连续对称性自发破缺；然而原子核是有限介观粒子体系，连续转动对称性 $[\hat{H}, \hat{\mathbf{J}}] = 0$ 在实验室系中是不可动摇的严格守恒律。微观平均场中出现的所谓“自发对称破缺”，仅仅是变分引入的内禀单粒子波包的内部属性；通过角动量投影技术，可从微观破缺波包中提取出具有确定良好自旋量子数的试探空间态矢（good-$I$ 态）；而在该受限物理子空间中进一步通过哈密顿量 Ritz 变分（Hill-Wheeler 方程）求解，便在实验室参考系中求得满足守恒律的物理近似本征态（详见第2.5节）。

\subsection{章节脉络与后续动力学演化接口}\label{sec:3-1-5}

本节所澄清的几何取向、内禀振动与量子涨落之间的物理边界，为全书第3章展开深入的低能动力学分析提供了清晰的逻辑阶梯：
\begin{itemize}
\item \textbf{第 3.2 节（稳定转轴与摆动的局域极限）}：当四极形状势能变硬、内禀形状振动被有效冻结时，我们将坐标系完全固定于三轴椭球主轴系，深入探究经典欧拉陀螺的运动稳定性，并在局域大自旋展开下精确推导纯转子的三维摆动声子激发频率及其两种轴对称极限；
\item \textbf{第 3.3 节（准粒子顺排与横/纵向摆动分支）}：将微观壳模型费米面与未配对价核子带入集体核芯，研究高 $j$ 准粒子轨道角动量顺排如何打破主轴旋转对称性，催生出具有丰富微观指纹的横向摆动（TW）与纵向摆动（LW）分支；
\item \textbf{第 3.5 节（手征几何、隧穿与有限核态）}：当三轴非轴对称集体核芯与一个高 $j$ 价粒子和一个高 $j$ 价空穴实现非共面角动量矢量耦合时（标量三重积 $\mathbf{j}_p \cdot (\mathbf{j}_h \times \mathbf{R}) \ne 0$，并不要求三矢量严格 $90^\circ$ 正交），体系在旋转参考系诱发自发手征对称破缺，进而在有限量子隧穿作用下于实验室系形成近简并的手征双重带（Chiral doublet bands）；
\item \textbf{第 3.6 节（$\gamma$ 形状模与取向模的耦合）}：重新放宽形状冻结假设，系统探讨内禀 $\gamma$ 形状振动与欧拉三维取向运动之间的非线性动力学杂化，阐明同自旋、同宇称能带在有限能隙下如何通过非对角矩阵元诱发强烈的态混合。
\end{itemize}
